\paragraph{Separating the repair objective from the search engine.}
For the remaining 5487/pin~24 discrepancy, a new default-off serial-router filter forbids occupied extension nodes and re-entry into retained tree nodes while preserving native retained seeds. We independently switch the native mixed delay/congestion objective to physical delay, retaining true switch delay and final STA. The new default control exactly reproduces all non-time fields of the previous native trace. Mixed cost with either soft congestion or hard capacity keeps the old timing plateau. Pure delay with soft congestion returns an overused path and rolls back. Pure delay with hard capacity reproduces the complete physical tree of the capacity-oracle candidate, improving the same two endpoint components while CPD remains unchanged. Thus the observed small benefit requires the tested objective/constraint combination; it is not unique to the external search engine. Four native analyses make six actual connection-search attempts. The three complete candidates reuse already certified FULL/STA only after exact tree/input/constraint matching. Incremental orchestration takes 68.542 s; the isolated build takes 52.104 s. These are development measurements on the zero-RC case, not a general native performance claim.

\paragraph{Paired successors do not require the earlier plateau steps.}
We give the strict state before the four endpoint-only steps and the endpoint state after them the same current-STA joint-repair opportunity. The current worst connections remain 76/pin~1 and 48/pin~1. Their capacity-constrained distances have no gain, but each relaxed fixed-tree distance improves by 0.185825799 ns. Only two new relaxed queries are needed: identical unconstrained inputs in the second state reuse the first answers. Compatible target paths identify the same five blocking nets, each with one affected sink. Both completed transactions pass native FULL/STA and reach 5.832274663220005 ns, improving eight setup components and regressing one relative to their own initial states. Therefore the four preceding endpoint steps are not necessary for this CPD reduction. Their five endpoint improvements remain in the later candidate, but do not enable a uniquely lower CPD in this experiment.

Strict CPD accepts the successor, improving the earlier 5.854659868 ns result by approximately 0.38235\%, and the original strong R1 6.018101129 ns result by 3.087792\%. The endpoint-Pareto branch rejects and retains its own state because \texttt{Mp[12].D[0]} setup slack worsens by 0.371651154 ns. The directly connected net 79/pin~1 delay increases by 0.379290599 ns; the difference shows that connection delay alone does not equal endpoint slack change. The accepted strict state's worst group migrates to six connections, including 5490/pin~5 and 6591/pin~4. Its physical resources are 45 SB pairs, 975 CB sources and planar wirelength 218190. The candidate after the plateau steps uses 48, 976 and 218274, but is not committed by the endpoint policy. No acceptance tolerance is introduced.

\paragraph{Native recovery reproduces the new CPD with fewer crossing resources.}
Two actual native sequential controls start from the common strict input and receive exactly the same two-target/five-blocker released scope, without preset target paths. Mixed-cost native recovery and physical-delay hard-capacity recovery both reach 5.832274663 ns and exactly the external candidate's complete setup/hold vectors. Both pass fresh FULL/STA; each strict branch accepts and each endpoint-Pareto comparison rejects. Mixed-cost native recovery uses 41 SB pairs and 971 CB sources, four fewer of each than the external candidate, with planar wirelength 218225, 35 more. The physical hard-capacity arm uses 45, 975 and 218190. Their seven-connection transactions take 0.354 and 0.307 s internally and require thirteen and seven native searches respectively. Complete control orchestration takes 170.627 s. This supports prioritizing native recovery, current-STA selection, released scope and repair opportunities; the released scope still comes from the automatic path/conflict analysis, so autonomous native scope discovery remains to be isolated.

The paired potential stage takes 50.061 s and fresh paired FULL/STA orchestration 80.998 s. Recorded transaction construction/audit timers are 37.212 and 37.991 s, but omit pre-registration/import work; the independently launched adaptive phases lack one enclosing end-to-end wall timer. Their overlapping times must not be summed into a claimed whole-flow speed. The next unified controller must time every child from process creation, including failures and rollback. No historical U timing is substituted. All quality results here remain local development evidence rather than a cross-circuit promotion.
